\subsection{DERIVATIVES OF ORDER n}

Let \(\mathbf{f}\) be a vector function of a real variable, defined, continuous and differentiable on an interval I. If the derivative \(\mathbf{f}'\) exists on a neighbourhood (with respect to I) of a point \(x_0 \in \mathrm{I}\) , and is differentiable at the point \(x_0\) , then its derivative is called the second derivative of \(\mathbf{f}\) at the point \(x_0\) , and is denoted by \(\mathbf{f}''(x_0)\) or \(\mathrm{D}^2\mathbf{f}(x_0)\) . If this second derivative exists at every point of I (which implies that \(\mathbf{f}'\) exists and is continuous on I), then \(x \mapsto \mathbf{f}''(x)\) is a vector function which one denotes by \(\mathbf{f}''\) or \(\mathrm{D}^2\mathbf{f}\) . We define, in the same way, recursively, the \(n^{th}\) derivative (or derivative of order \(n\) ) of \(\mathbf{f}\) , and denote it by \(\mathbf{f}^{(n)}\) or \(\mathrm{D}^n\mathbf{f}\) ; by definition, its value at the point \(x_0 \in \mathrm{I}\) is the derivative of the function \(\mathbf{f}^{(n-1)}\) at the point \(x_0\) : this definition presupposes the existence of all the derivatives \(\mathbf{f}^{(k)}\) of order \(k \leqslant n - 1\) on a neighbourhood of \(x_0\) relative to I, and the differentiability of \(\mathbf{f}^{(n-1)}\) at the point \(x_0\) .

We will say that \(\mathbf{f}\) is \(n\) times differentiable at the point \(x_0\) (resp. in an interval) if it admits an \(n^{th}\) derivative at this point (resp. in this interval). One says that \(\mathbf{f}\) is indefinitely differentiable on I if for each integer \(n > 0\) it admits a derivative of order \(n\) on I.

By induction on \(m\) one sees that

\[
\mathbf {D} ^ {m} (D ^ {n} \mathbf {f}) = \mathbf {D} ^ {m + n} \mathbf {f}.\tag{1}
\]

More precisely, when one of the two terms in (1) is defined, then so is the other, and is equal to it.

PROPOSITION 1. The set of vector functions defined on an interval \(I \subset R\) , taking values in a given topological vector space E, and having an \(n^{th}\) derivative on I, is a vector space over R, and \(f \mapsto D^{n}f\) is a linear mapping of this space into the vector space of linear mappings from I into E.

One proves the formulae

\[
\mathrm{D} ^ {n} (\mathbf {f} + \mathbf {g}) = \mathrm{D} ^ {n} \mathbf {f} + \mathrm{D} ^ {n} \mathbf {g}\tag{2}
\]

\[
\mathrm{D} ^ {n} (\mathbf {f} a) = \mathrm{D} ^ {n} \mathbf {f}. a\tag{3}
\]

by induction on n when f and g have an \(n^{th}\) derivative on I (a being constant).

PROPOSITION 2 (``Leibniz' formula''). Let E, F, G be three topological vector spaces over R, and \((\mathbf{x}, \mathbf{y}) \mapsto [\mathbf{x}.\mathbf{y}]\) a continuous bilinear mapping of \(E \times F\) into G. If f (resp. g) is defined on an interval \(I \subset R\) , takes its values in E (resp. F) and has an \(n^{th}\) derivative on I, then {[}f.g{]} has an \(n^{th}\) derivative on I, given by the formula

\[
\mathrm{D} ^ {n} [ \mathbf {f}. \mathbf {g} ] = [ \mathbf {f} ^ {(n)}. \mathbf {g} ] + \binom {n} {1} [ \mathbf {f} ^ {(n - 1)}. \mathbf {g} ^ {\prime} ] + \dots + \binom {n} {p} [ \mathbf {f} ^ {(n - p)}. \mathbf {g} ^ {(p)} ] + \dots + [ \mathbf {f}. \mathbf {g} ^ {(n)} ].\tag{4}
\]

Formula (4) is proved by induction on \(n\) (using the relation \(\binom{n}{p} = \binom{n-1}{p} + \binom{n-1}{p-1}\) for the binomial coefficients).

In the same way one can verify the following formula (where the hypotheses are the same as in prop. 2):

\[
\begin{array}{c} [ \mathbf {f} ^ {(n)}. \mathbf {g} ] + (- 1) ^ {n - 1} [ \mathbf {f}. \mathbf {g} ^ {(n)} ] = D \big ([ \mathbf {f} ^ {(n - 1)}. \mathbf {g} ] - [ \mathbf {f} ^ {(n - 2)}. \mathbf {g} ^ {\prime} ] + \dots \\ + (- 1) ^ {n - 1} [ \mathbf {f}. \mathbf {g} ^ {(n - 1)} ] \big). \end{array}\tag{5}
\]

The preceding propositions have been stated for functions that are n times differentiable on an interval; we leave it to the reader to formulate the analogous propositions for functions that are n times differentiable at a point.

